\section{JUSTIFICACIÓN}

La justificación del presente proyecto expone las razones que sustentan su desarrollo, articulando la necesidad identificada en el taller de mecanizado, el beneficio técnico-operativo de la solución mecatrónica, la importancia de la innovación tecnológica, el aporte al conocimiento académico y los beneficiarios directos e indirectos, estructurándose en tres dimensiones fundamentales:

\subsection{Justificación técnica}
Desde la perspectiva técnica, la problemática de los acabados superficiales deficientes y la pérdida de precisión dimensional en el mecanizado radica en la imposibilidad de evaluar de forma continua el desgaste de los insertos durante el corte. En este contexto, el desarrollo de un sistema periférico no invasivo representa una solución idónea frente a la arquitectura cerrada de los tornos CNC de MAINSA, permitiendo instrumentar la máquina sin intervenir el cableado interno ni alterar la estabilidad del equipo. La integración sinérgica de sensores de corriente en el husillo y acelerometría en la torreta captura tanto la variación en los esfuerzos globales de corte como las oscilaciones dinámicas de alta frecuencia originadas por la degradación del filo.

Asimismo, la incorporación de procesamiento inteligente en el borde (\textit{Edge AI}) sobre una unidad embebida resuelve la necesidad de respuesta inmediata ante anomalías operativas de mecanizado. Al realizar la inferencia algorítmica de manera local, se prescinde de la transmisión masiva de señales hacia servidores remotos, eliminando la latencia de red y asegurando una supervisión continua y robusta. Esta arquitectura técnica garantiza que el sistema opere de forma confiable frente al elevado ruido electromagnético característico del taller industrial, proporcionando un diagnóstico oportuno del estado del inserto antes de que se produzcan mermas en las piezas trabajadas.

\subsection{Justificación económica}
En el ámbito económico, el desarrollo de una solución tecnológica a nivel nacional en Bolivia ofrece una viabilidad significativamente superior en comparación con la adquisición de sistemas comerciales de importación. El desarrollo local prescinde de los elevados costos asociados a la logística de transporte internacional, trámites aduaneros, pago de aranceles e intermediación comercial, haciendo accesible una tecnología de monitoreo avanzada para talleres metalmecánicos medianos como MAINSA. Esta accesibilidad económica viabiliza la modernización tecnológica de maquinaria heredada sin incurrir en inversiones de capital prohibitivas para el entorno productivo local.

Adicionalmente, el beneficio económico del sistema se manifiesta directamente en la reducción de costos operativos dentro del taller. El desgaste inadvertido de las herramientas suele provocar el descarte definitivo de piezas semielaboradas o demandar costosas jornadas de retrabajo para corregir tolerancias y rugosidades fuera de norma, además de consumir horas improductivas de máquina por paradas imprevistas. Al proporcionar una alerta temprana sobre la pérdida de filo, el sistema permite sustituir los insertos en el momento técnicamente oportuno, maximizando la vida útil de cada filo de carburo, disminuyendo el desperdicio de materia prima y optimizando la rentabilidad global del proceso de manufactura.

\subsection{Justificación académica y social}
En el plano académico, esta investigación aporta un caso de aplicación real y metodológicamente validado para la carrera de Ingeniería Mecatrónica de la Universidad Católica Boliviana, articulando sinérgicamente cuatro pilares formativos: mecánica de manufactura CNC, instrumentación y acondicionamiento analógico de señales, procesamiento digital embebido con inteligencia artificial (\textit{TinyML}) y telemática industrial (IoT). El proyecto genera metodologías de adquisición de datos en planta y modelos de desgaste reproducibles que servirán como base académica y experimental para futuras investigaciones de grado en mantenimiento predictivo y manufactura inteligente.

En el plano social, el proyecto genera un impacto positivo directo en las condiciones laborales de los operarios de tornos CNC en MAINSA. El sistema mitiga sustancialmente la exposición a riesgos de accidentes laborales por desprendimiento violento de esquirlas o fragmentos incandescentes ante roturas catastróficas de insertos, al tiempo que reduce el estrés visual y la fatiga asociados a la supervisión manual continua en caliente. Finalmente, como beneficio indirecto, los clientes del taller metalmecánico reciben componentes maquinados con estricta adherencia a tolerancias dimensionales y rugosidades superficiales, fortaleciendo la competitividad y confiabilidad de la industria manufacturera regional.
